\documentclass[12pt,a4paper]{article}
\usepackage[utf8]{inputenc}
\usepackage{amsmath}
\usepackage{amsfonts}
\usepackage{graphicx}
\usepackage{hyperref}
\usepackage{geometry}
\usepackage{lmodern}
\geometry{a4paper, margin=1in}

\title{\textbf{Mathematical Modeling of Epidemic Spread Using the SIR Model}}
\author{Milica Nikolic}
\date{August 2026}

\begin{document}

\maketitle

\begin{abstract}
TODO
\end{abstract}

\section{Introduction}
TODO 

\section{The SIR Model}

The Susceptible-Infectious-Recovered (SIR) model stands as one of the most fundamental and influential frameworks in
mathematical epidemiology, serving as the cornerstone for understanding infectious disease dynamics across diverse
populations and contexts.

The model is based on dividing the host population (humans, in the case of this project) into a small number of compartments, each containing individuals that are identical in terms of their status with respect to the disease in question. In the SIR model, there are three primary compartments:
\begin{itemize}
    \item Susceptible ($S$): Individuals who have no immunity to the infectious agent, so they might become infected if exposed.
    \item Infectious ($I$): Individuals who are currently infected and can transmit the infection to susceptible individuals whom they contact.
    \item Removed ($R$): Individuals who are immune to the infection (either through recovery or death), and consequently do not affect the transmission dynamics in any way when they contact other individuals.
\end{itemize}
It is traditional to denote the number of individuals in each of these compartments as $S$, $I$, and $R$, respectively. The total host population size is defined as $N = S + I + R$.

Having compartmentalized the population, the next step is to define the mathematical rules that govern how individuals transition between these states over time. These transitions are expressed as a system of non-linear ordinary differential equations (ODEs) that capture the specific rates of infection and recovery, allowing us to simulate the continuous spread and eventual decline of the disease.

\subsection{Basic Assumptions}
The classic SIR model relies on several key simplifying assumptions:
\begin{itemize}
    \item \textbf{Constant Population:} The total population $N$ remains constant over time ($N = S + I + R$). Births, migrations, and deaths by other causes are ignored.
    \item \textbf{Homogeneous Mixing:} It is assumed that every individual has an equal probability of coming into contact with any other individual in the population.
    \item \textbf{Permanent Immunity:} Once an individual recovers, they acquire lifelong immunity and cannot be reinfected.
\end{itemize}

\subsection{Model Parameters}
The flow of individuals between the compartments is governed by two primary rates:
\begin{itemize}
    \item \textbf{$\beta$ }(Transmission Rate): This is the 'per capita' rate at which the infection is transmitted from an infectious individual to a susceptible one. In broader epidemiological terms, $\beta$ is usually understood as the product of two distinct factors: the average number of contacts a person makes per unit of time, and the probability that a disease is actually transmitted during one of those contacts.
    \item \textbf{$\gamma$ }(Recovery Rate): This parameter represents the rate at which infected individuals recover (or are otherwise removed from the infectious pool, such as through isolation or death). Its reciprocal, $1/\gamma$, represents the mean infectious period—the average amount of time an individual remains capable of transmitting the disease before moving to the removed compartment.
    \item \textbf{$N$} (Total Population Size): The total number of individuals in the host population, defined as $N = S + I + R$. As your text notes, for the differential equations to treat $S$, $I$, and $R$ as smooth continuous variables rather than discrete integers, $N$ must be sufficiently large.
    \item \textbf{$R_0$ }(Basic Reproduction Number): Defined in your text as $R_0 = \frac{\beta N}{\gamma}$, this is the most critical threshold quantity in epidemic modeling. It represents the expected number of secondary infections generated by a single infected individual introduced into a completely susceptible population ($S(0) = N$).
    \begin{itemize}
        \item If $R_0 > 1$, each infection spawns more than one new infection, meaning the disease will spread and an epidemic will occur.
        \item If $R_0 < 1$, the number of cases will decrease over time, and the outbreak will naturally die out.
    \end{itemize}
    
\end{itemize}

\subsection{System of Differential Equations}
The dynamic transitions between the Susceptible ($S$), Infectious ($I$), and Recovered ($R$) states are mathematically described by the following system of non-linear ordinary differential equations:


\begin{align}
    \frac{dS}{dt} &= -\frac{\beta S I}{N} \label{eq:sir_S} \\
    \frac{dI}{dt} &= \frac{\beta S I}{N} - \gamma I \label{eq:sir_I} \\
    \frac{dR}{dt} &= \gamma I \label{eq:sir_R}
\end{align}

Let us now explain the equations presented above.

The first equation \eqref{eq:sir_S} represents the rate of change for the susceptible population ($S$). Intuitively, this number decreases over time as individuals contract the disease and leave this group, which accounts for the negative sign. The multiplication of $S$ and $I$ stems from the assumption that the population mixes randomly. It reflects the fact that the likelihood of a disease transmission event depends directly on a susceptible individual coming into contact with an infected one. We divide by the total population size, $N$, because the term $\frac{I}{N}$ represents the fraction of the population that is currently infected. This fraction gives the realistic probability that a susceptible person will encounter someone who is sick. By dividing by $N$, we make sure the model assumes people have a normal amount of daily contacts, rather than incorrectly assuming they meet way more people just because they live in a larger population. Finally, the transmission parameter, $\beta$, dictates the overall rate at which these infectious contacts occur.


As the number of susceptible individuals decreases, they must transition into the infected group ($I$). This inflow of newly infected individuals explains the first term on the right-hand side of the second equation \eqref{eq:sir_I}. At the same time, infected individuals are recovering at a specific rate, $\gamma$. This continuous outflow from the infected group accounts for the subtracted second term.

The third equation \eqref{eq:sir_R} describes the change in the removed or recovered population ($R$). It is intuitive that this number---which, as a reminder, encompasses both successful recoveries and disease-induced fatalities---will grow over time. The rate at which this group increases is exactly equal to the rate at which individuals exit the infected group.

\section{Numerical Methods}
It is important to note that the system of ordinary differential equations governing the SIR model is nonlinear due to the interaction between the susceptible and infected populations. Because of this nonlinearity, these ODEs do not have a simple analytical solution. Therefore, to analyze the dynamics of the epidemic, it is necessary to employ numerical methods to approximate the solutions to these equations.

\subsection{Euler Method}
The Euler method is a fundamental first-order numerical procedure used for approximating the solutions of ordinary differential equations (ODEs) with a given initial value. It calculates the solution curve step-by-step by using the derivative (slope) of the function at a known point to estimate the function's value at a subsequent point. 

Given an initial value problem of the form $\frac{dy}{dt} = f(t, y)$ with the initial condition $y(t_0) = y_0$, the Euler method discretizes the time domain into steps of size $\Delta t$. The value of the state variable at the next time step, $y_{n+1}$, is computed using the following explicit formula:

\begin{equation}
    y_{n+1} = y_n + \Delta t \cdot f(t_n, y_n) \label{eq:euler_general}
\end{equation}

While the Euler method is computationally straightforward and easy to implement, it is a first-order method. Consequently, its accuracy and numerical stability are highly dependent on selecting a sufficiently small time step, $\Delta t$, to minimize the local truncation error at each iteration.

\subsection{Runge-Kutta Method (RK4)}

The classic fourth-order Runge-Kutta method (often abbreviated as RK4) is a highly accurate and widely used numerical technique for solving ordinary differential equations. While the first-order Euler method relies solely on the derivative at the beginning of a time step, RK4 achieves greater precision by calculating four distinct slopes within each interval and computing their weighted average.

For a given initial value problem $\frac{dy}{dt} = f(t, y)$, the method determines the next value, $y_{n+1}$, using the following sequence of calculations:

\begin{align}
    k_1 &= f(t_n, y_n) \label{eq:rk4_k1} \\
    k_2 &= f\left(t_n + \frac{\Delta t}{2}, y_n + \frac{\Delta t}{2} k_1\right) \label{eq:rk4_k2} \\
    k_3 &= f\left(t_n + \frac{\Delta t}{2}, y_n + \frac{\Delta t}{2} k_2\right) \label{eq:rk4_k3} \\
    k_4 &= f(t_n + \Delta t, y_n + \Delta t k_3) \label{eq:rk4_k4}
\end{align}

Once these four slopes are evaluated, the state at the next time step is computed as:

\begin{equation}
    y_{n+1} = y_n + \frac{\Delta t}{6} (k_1 + 2k_2 + 2k_3 + k_4) \label{eq:rk4_final}
\end{equation}

In this weighted average, the midpoint slopes ($k_2$ and $k_3$) are given twice the weight of the endpoints ($k_1$ and $k_4$). Because RK4 is a fourth-order method, the global truncation error is on the order of $O(\Delta t^4)$. This makes it vastly more accurate and stable than the Euler method, allowing for larger time steps without sacrificing precision.

\section{Results and Analysis}

\subsection{Baseline Simulation}
% To insert a generated plot later, uncomment the lines below and place the image in the same folder.
% \begin{figure}[h]
%     \centering
%     % \includegraphics[width=0.75\textwidth]{sir_baseline_plot.png}
%     \caption{Epidemic curve showing S, I, and R populations over time (Baseline parameters).}
%     \label{fig:baseline}
% \end{figure}

\subsection{Impact of Varying Parameters}
TODO

\section{Conclusion}
TODO

\end{document}
